\documentclass[11pt]{article}
% Working draft. Swap \documentclass for a venue style (e.g. iclr/neurips) when targeting one.
\usepackage[margin=1in]{geometry}
\usepackage{times}
\usepackage{hyperref}
\usepackage{url}
\usepackage{graphicx}
\graphicspath{{figures/}}
\usepackage{booktabs}
\usepackage{amsmath}
\usepackage{amssymb}
\usepackage{multirow}
\usepackage[numbers]{natbib}
\usepackage[capitalize,noabbrev]{cleveref}
\usepackage{xspace}

% Minimal math/notation macros for the identity-trap manuscript.
\newcommand{\BA}{\mathrm{BA}}                 % balanced accuracy
\newcommand{\subjf}{s_{\mathrm{subj}}}        % subject fraction (LEACE)
\newcommand{\reve}{REVE\xspace}
\newcommand{\tcm}{TCM\xspace}
\newcommand{\fmscope}{\textsc{FMScope}\xspace}

% AUTO-GENERATED by paper/generate.py -- do not edit by hand.
\newcommand{\fmmiPoolTrap}{17}
\newcommand{\fmmiPoolN}{17}
\newcommand{\fmmiPtTrap}{0}
\newcommand{\fmmiPtN}{12}
\newcommand{\tcmmiPoolLift}{7}
\newcommand{\tcmmiN}{13}
\newcommand{\tcmmiPtTrap}{0}
\newcommand{\fmerpPoolTrap}{14}
\newcommand{\fmerpPoolN}{22}
\newcommand{\fmerpPtTrap}{0}
\newcommand{\fmerpPtN}{23}
\newcommand{\tcmerpPoolLift}{7}
\newcommand{\tcmerpN}{22}
\newcommand{\tcmerpPtTrap}{0}
\newcommand{\fmssvepPoolTrap}{3}
\newcommand{\fmssvepPoolN}{6}
\newcommand{\fmssvepPtTrap}{1}
\newcommand{\fmssvepPtN}{6}
\newcommand{\tcmssvepPoolLift}{1}
\newcommand{\tcmssvepN}{6}
\newcommand{\tcmssvepPtTrap}{0}
\newcommand{\fmPtTrapTotal}{1}
\newcommand{\fmPtNTotal}{41}
\newcommand{\fmPoolTrapTotal}{34}
\newcommand{\reveSubjfmiLo}{0.85}
\newcommand{\reveSubjfmiHi}{0.98}
\newcommand{\reveSubjfmiMean}{0.92}
\newcommand{\reveSubjferpLo}{0.55}
\newcommand{\reveSubjferpHi}{0.95}
\newcommand{\reveSubjferpMean}{0.86}
\newcommand{\reveSubjfssvepLo}{0.67}
\newcommand{\reveSubjfssvepHi}{0.90}
\newcommand{\reveSubjfssvepMean}{0.81}
\newcommand{\bnciFmPoolRaw}{0.667}
\newcommand{\bnciFmPoolFree}{0.963}
\newcommand{\bnciFmPtRaw}{0.537}
\newcommand{\bnciFmPtFree}{0.539}
\newcommand{\bnciTcmPoolRaw}{0.537}
\newcommand{\bnciTcmPoolFree}{0.648}
\newcommand{\bnciTcmPtRaw}{0.513}
\newcommand{\bnciTcmPtFree}{0.509}
\newcommand{\nakRaw}{0.877}
\newcommand{\nakFree}{0.899}
\newcommand{\nakNsub}{9}
\newcommand{\nakNtr}{270}
\newcommand{\nakTcmPtRaw}{0.840}
\newcommand{\nakTcmPtFree}{0.814}
 % auto-generated macros (paper/generate.py): every headline number
% AUTO-GENERATED by paper/generate.py.
\newcommand{\trapCountTable}{\begin{tabular}{lcccc}
\toprule
Paradigm & REVE pooled & REVE per-trial & TCM pooled-lift & TCM per-trial \\
\midrule
MI & 17/17 & 0/12 & 7/13 & 0/13 \\
ERP & 14/22 & 0/23 & 7/22 & 0/22 \\
SSVEP & 3/6 & 1/6 & 1/6 & 0/6 \\
\bottomrule
\end{tabular}}
\newcommand{\reproTable}{\begin{tabular}{lcccc}
\toprule
Dataset & REVE pooled & REVE per-trial & TCM pooled & TCM per-trial \\
\midrule
BNCI2014-001 & 0.667$\to$0.963 & 0.537$\to$0.539 & 0.537$\to$0.648 & 0.513$\to$0.509 \\
Schirrmeister2017 & 0.690$\to$0.964 & 0.546$\to$0.551 & 0.821$\to$1.000 & 0.592$\to$0.591 \\
Stieger2021 & 0.981$\to$1.000 & 0.730$\to$0.710 & 0.989$\to$1.000 & 0.790$\to$0.745 \\
\bottomrule
\end{tabular}}
\newcommand{\headToHeadTable}{\begin{tabular}{lcccc}
\toprule
ERP dataset & REVE pt raw & TCM pt raw & REVE subj.\ frac & TCM subj.\ BA \\
\midrule
Huebner2017 & 0.648 & 0.813 & 0.92 & 0.20 \\
ErpCore2021-ERN & 0.750 & 0.797 & 0.68 & 0.29 \\
Lee2019-ERP & 0.675 & 0.760 & 0.78 & 0.12 \\
BrainInvaders2012 & 0.628 & 0.690 & 0.93 & 0.10 \\
ErpCore2021-N170 & 0.628 & 0.686 & 0.90 & 0.58 \\
Huebner2018 & 0.635 & 0.549 & 0.89 & 0.12 \\
BrainInvaders2014a & 0.575 & 0.505 & 0.95 & 0.17 \\
\bottomrule
\end{tabular}}
 % auto-generated tables

\title{The Identity Trap Is a Pooling Artifact:\\
A Per-Trial, Cross-Paradigm, Cross-Method Audit of EEG Decoders}
\author{Morgan Hough\\\texttt{NeuroTechX}}
\date{\today}

\begin{document}
\sloppy  % tolerate looser inter-word spacing to avoid minor overfull boxes
\maketitle

\begin{abstract}
\fmscope{} \citep{fmscope} reports that EEG foundation models (FMs) suffer an
\emph{identity trap}: a subject-axis erasure (LEACE; \citealp{belrose2023leace}) that removes
subject identity from a frozen FM's features \emph{raises} downstream task accuracy, implying
the apparent skill rode on subject leakage. We show this effect is governed by the
\emph{evaluation metric}, not by foundation models. The published erasure pools each
$(\text{subject},\text{class})$ into a single recording whose predictions are mean-pooled; that
pooling crushes within-recording variance and manufactures a large recovered-skill lift. We
re-run the audit two ways---the published \emph{pooled} erasure and an \emph{honest per-trial}
erasure (each window its own recording, grouped by subject, $n\gg p$)---across the entire MOABB
benchmark spanning motor imagery (MI), P300 (ERP), and SSVEP, for two feature families: the
\reve{} foundation model and a classical Wang Topographic Component Model (\tcm{};
\citealp{wang2000tcm}) control with far weaker identity encoding. Under the pooled metric the
trap reproduces faithfully---\reve{} traps on \fmPoolTrapTotal{} datasets, including every MI
dataset (\fmmiPoolTrap{}/\fmmiPoolN{})---and the canonical case BNCI2014-001 (the
\fmscope{} paper's own MI trap) recovers from $\bnciFmPoolRaw{}\rightarrow\bnciFmPoolFree{}$. Under
the per-trial metric the trap essentially vanishes: \fmPtTrapTotal{}/\fmPtNTotal{} \reve{} datasets
trap. Crucially the classical \tcm{} control reproduces the \emph{same} pooled$\to$per-trial
collapse despite encoding far less identity, which localizes the trap to the pooled metric rather
than to foundation models. \reve{}'s representation is heavily subject-dominated (subject fraction
$\reveSubjfmiMean{}$/$\reveSubjferpMean{}$/$\reveSubjfssvepMean{}$ for MI/ERP/SSVEP) yet shows no
robust per-trial trap, and on several clean ERP contrasts the plain-SVD control \emph{out-decodes}
the foundation model per trial while carrying a fraction of the identity. We conclude the identity
trap, as published, is a pooling artifact; the honest single-trial test is the one that should be
reported.
\end{abstract}

\section{Introduction}
EEG foundation models promise transferable, subject-agnostic representations, but they are trained
and evaluated on small-$N$, subject-rich corpora where any apparent task skill risks riding on
subject identity. \fmscope{} \citep{fmscope} formalizes this worry with a \emph{subject-axis
erasure}: fit a concept-erasure transform (LEACE; \citealp{belrose2023leace}) that linearly removes
subject identity from a frozen model's features, then re-decode the task. If accuracy \emph{rises}
after erasure---an ``identity-free lift''---the masked subject signal was hurting, and the model was
trapped into encoding identity rather than task. The paper reports this trap for several EEG FMs.

We ask a narrower, sharper question: \emph{is the trap a property of foundation models, or of the
erasure's evaluation protocol?} The published protocol scores erasure at the level of
\emph{recordings}, where one recording aggregates all trials of a given $(\text{subject},
\text{class})$ and its per-trial predictions are mean-pooled into a single label. With only a
handful of subjects per class, this pooling collapses the sample to a few high-variance points, so
a linear probe's balanced accuracy swings wildly and a post-erasure ``lift'' is easy to manufacture.
We contrast this with an \emph{honest per-trial} erasure in which every window is its own recording,
grouped by subject so train/test never share a subject, giving $n\gg p$ and a stable estimate.

\paragraph{Contributions.}
(i) We run both the pooled (published) and per-trial erasure across the full MOABB
\citep{jayaram2018moabb} benchmark---MI, ERP (P300), and SSVEP---for the \reve{} foundation model.
(ii) We add a classical \emph{control}, the Wang Topographic Component Model (\tcm{};
\citealp{wang2000tcm}), an unsupervised per-mode trilinear decomposition whose loadings encode far
less subject identity than an FM. (iii) We show the trap reproduces under pooling for \emph{both}
the FM and the classical control, and dissolves under the per-trial metric for both---evidence that
the trap is metric-induced. (iv) We report a decoder--identity dissociation: \reve{} is more
identity-dominated yet not more trap-prone per trial, and a plain SVD control out-decodes it per
trial on several clean ERP contrasts.

\section{Related work}
\paragraph{Identity confounds in EEG decoding.} EEG is notoriously subject-specific, and decoders
that look skilled under subject-mixed evaluation often collapse across subjects; MOABB
\citep{jayaram2018moabb} standardises within- vs.\ cross-subject splits precisely to expose this.
\fmscope{} \citep{fmscope} sharpens the concern for foundation models with a subject-axis erasure,
reporting that frozen EEG-FM features are ``identity-trapped.'' We do not dispute the measurement;
we show its verdict is governed by how erasure is scored.

\paragraph{Linear concept erasure.} Removing a concept (here, subject identity) from a representation
by a linear map is well studied: iterative nullspace projection \citep{ravfogel2020inlp}, adversarial
relaxations \citep{ravfogel2022rlace}, and the closed-form LEACE \citep{belrose2023leace} used here
and by \fmscope{}. These methods erase a direction; they say nothing about how the downstream metric
is aggregated, which is where the trap originates.

\paragraph{Pooling and pseudoreplication.} Averaging correlated within-unit observations into one
point per unit before testing is the classical pseudoreplication error \citep{hurlbert1984}: it
shrinks within-unit variance and inflates apparent effects. The pooled erasure
(\cref{eq:pool}) is exactly this---one mean-pooled point per $(\text{subject},\text{class})$---so the
``trap'' is the statistical signature of pooling, recovered here for a foundation model and an
unsupervised classical control alike.

\section{Methods}
\paragraph{Subject-axis erasure.} For features $X\in\mathbb{R}^{n\times d}$ with subject labels and
task labels, we fit LEACE to linearly erase the subject axis and report task balanced accuracy
($\BA$) before (\emph{raw}) and after (\emph{identity-free}) erasure, plus a subject-probe $\BA$ and
the LEACE subject fraction $\subjf$ (the share of representational variance attributable to subject
identity). A row is \emph{interpretable} only if the raw task decode clears a $0.55$ gate; among
interpretable rows we call a \emph{trap} when the identity-free $\BA$ exceeds the raw $\BA$ by more
than $0.02$, \emph{task-carried} when it does not, and \emph{no-transfer} below the gate.

\paragraph{Pooled vs.\ per-trial.} Let the data be $S$ subjects $\times$ $C$ classes, with $m_{sc}$
trials in cell $(s,c)$ and $n=\sum_{s,c} m_{sc}$ trials total. A linear probe $w$ gives a per-trial
score $\sigma(w^{\top}\phi_i)$. The \emph{pooled} protocol (as published) forms one recording per
$(s,c)$ whose label is the \emph{mean} of its per-trial scores,
\begin{equation}
\hat p_{sc} \;=\; \frac{1}{m_{sc}}\sum_{i\in(s,c)} \sigma\!\left(w^{\top}\phi_i\right),
\label{eq:pool}
\end{equation}
and scores $\BA$ over the $SC$ aggregated points. Mean-pooling shrinks each cell's score variance by
$\approx 1/m_{sc}$, so the $SC$ points are far more separable than individual trials; with $SC$
small (a handful of subjects $\times$ two classes), \emph{both} the fitted LEACE erasure direction
and the post-erasure $\BA$ are high-variance, and a spurious post-erasure lift arises easily. The
\emph{per-trial} protocol scores the same probe over all $n$ trials with stratified-$k$-fold CV
\emph{grouped by subject} (no subject split across folds), so the effective sample is $n\gg SC$ and
the estimate is stable. The trap is the gap between these two estimators, not a property of $\phi$.

\paragraph{Feature families.} \reve{} is a transformer EEG foundation model; we read its block-6
features under its published input contract (0.5--99.5\,Hz, 200\,Hz). The \tcm{} control is the
Wang/Begleiter/Porjesz per-mode trilinear decomposition: a spatial basis and a temporal basis are
obtained by separate variance-thresholded SVDs over all trials (unsupervised), and each trial's
compact loading $L=B^{\top}xC$ is the feature. Both families pass through the identical erasure.

\paragraph{Corpus and robustness.} We sweep every MOABB MI (\texttt{LeftRightImagery}), ERP
(\texttt{P300}, binary Target/NonTarget), and SSVEP ($n_{\text{classes}}=2$) dataset under a common
broadband contract. Giant cohorts are loaded under a 24-subject cap to bound memory
(\texttt{Lee2019\_ERP} otherwise OOMs at 54 subjects), and the broadband $f_{\max}$ is auto-clamped
below each cohort's native Nyquist (recovering 128\,Hz sets such as \texttt{BrainInvaders2012}).
Every number and table in this manuscript is tangled directly from the committed audit CSVs by
\texttt{paper/generate.py}.

\section{Results}
\paragraph{The trap is a pooling artifact.} \Cref{tab:trapcount} contrasts pooled and per-trial trap
counts. Under pooling \reve{} traps on $\fmmiPoolTrap{}/\fmmiPoolN{}$ MI, $\fmerpPoolTrap{}/\fmerpPoolN{}$
ERP, and $\fmssvepPoolTrap{}/\fmssvepPoolN{}$ SSVEP datasets ($\fmPoolTrapTotal{}$ total); under the
per-trial metric only $\fmPtTrapTotal{}/\fmPtNTotal{}$ trap. The classical \tcm{} control shows the same
collapse---pooled lifts on $\tcmmiPoolLift{}/\tcmmiN{}$ MI and $\tcmerpPoolLift{}/\tcmerpN{}$ ERP
datasets, zero per-trial traps (\cref{fig:trapfrac}). \Cref{fig:liftscatter} shows the mechanism
per dataset: pooled erasure lift is large and positive, while the matched per-trial lift collapses
into a band around zero---for the FM and the control alike.

\begin{table}[t]\centering
\caption{Identity-trap counts (datasets trapping / interpretable datasets) by paradigm, feature
family, and erasure metric. The trap is pervasive under pooling and near-absent per trial, for both
the foundation model and the classical control.}
\label{tab:trapcount}
\trapCountTable
\end{table}

\begin{figure}[t]\centering
\includegraphics[width=0.62\textwidth]{fig_trap_fraction.png}
\caption{Fraction of interpretable datasets that trap, by paradigm and erasure metric. Pooled
erasure (dark bars) traps most datasets---every MI set for \reve{}---while the per-trial metric
(light bars) is near-zero, for both the foundation model and the Wang TCM control.}
\label{fig:trapfrac}
\end{figure}

\begin{figure}[t]\centering
\includegraphics[width=0.52\textwidth]{fig_lift_scatter.png}
\caption{Per-dataset erasure lift, pooled (x) vs.\ per-trial (y). Each point is one dataset
(\reve{}=circle, Wang TCM=triangle; colour = paradigm). Points sit at large positive pooled lift
but inside the per-trial no-trap band ($|{\cdot}|<\liftEps$, grey), i.e.\ the lift is a property of
the pooled metric, not the model or the feature family.}
\label{fig:liftscatter}
\end{figure}

\paragraph{Canonical reproduction.} \Cref{tab:repro} reproduces the \fmscope{} paper's own
motor-imagery trap. BNCI2014-001 (the ds004362 class) recovers under \reve{} pooling from
$\bnciFmPoolRaw{}\rightarrow\bnciFmPoolFree{}$ but is flat per trial
($\bnciFmPtRaw{}\rightarrow\bnciFmPtFree{}$, below gate); the same dataset under the \tcm{} control
lifts pooled ($\bnciTcmPoolRaw{}\rightarrow\bnciTcmPoolFree{}$) and is flat per trial.

\begin{table}[t]\centering
\caption{Canonical MI trap reproduction. Pooled erasure recovers masked skill (the published
phenomenon) for both the FM and the control; per-trial erasure does not.}
\label{tab:repro}
\reproTable
\end{table}

\paragraph{Identity dominance without a per-trial trap.} \reve{}'s representation is heavily
subject-dominated---subject fraction in $[\reveSubjfmiLo{},\reveSubjfmiHi{}]$ (MI),
$[\reveSubjferpLo{},\reveSubjferpHi{}]$ (ERP), $[\reveSubjfssvepLo{},\reveSubjfssvepHi{}]$ (SSVEP)---yet
shows no robust per-trial trap. \Cref{tab:h2h} pairs per-trial decode with identity encoding on
representative ERP datasets: \reve{} carries far more identity than \tcm{} but is not more
trap-prone, and on clean contrasts (Huebner2017, ERPCore-ERN, Lee2019-ERP) the plain-SVD \tcm{}
control \emph{out-decodes} the foundation model per trial. \Cref{fig:identity} makes the
dissociation explicit across all per-trial datasets: \reve{}'s subject fraction is uniformly high,
yet it is uncorrelated with whether the task decodes per trial---identity dominance does not buy a
trap.

\begin{figure}[t]\centering
\includegraphics[width=0.62\textwidth]{fig_identity_vs_decode.png}
\caption{\reve{} per-trial: identity dominance (subject fraction, x) vs.\ per-trial task decode
(raw balanced accuracy, y), coloured by verdict. Identity fraction is high throughout but does not
predict the per-trial outcome---the representation is subject-dominated whether or not the task
survives erasure.}
\label{fig:identity}
\end{figure}

\begin{table}[t]\centering
\caption{Per-trial decode and identity encoding on representative ERP datasets. \reve{}'s subject
fraction is large, but the classical control decodes the task as well or better per trial with a
fraction of the identity. (\reve{} subject fraction and \tcm{} subject-probe BA are related but
distinct identity metrics; compare directions, not magnitudes.)}
\label{tab:h2h}
{\small\headToHeadTable}
\end{table}

\paragraph{The one exception is small-$N$ noise.} The single per-trial \reve{} trap across all
\fmPtNTotal{} datasets is \texttt{Nakanishi2015}---the smallest cohort ($\nakNsub{}$ subjects,
$\nakNtr{}$ windows), a $\nakRaw{}\rightarrow\nakFree{}$ lift barely over the $0.02$ threshold, and
\emph{not} reproduced by the \tcm{} control on the same data
($\nakTcmPtRaw{}\rightarrow\nakTcmPtFree{}$, a decrease). It is a borderline small-sample fluctuation,
not a robust counterexample.

\section{Discussion}
The identity trap, as published, is real \emph{under the pooled metric} and reproduces faithfully for
a foundation model across all three EEG paradigms. But it also reproduces for an unsupervised
classical decomposition that encodes far less subject identity, and it dissolves for both under an
honest per-trial erasure. A pathology that appears in a near-identity-free control and disappears when
the sample is not collapsed is a property of the estimator, not of the model: pooling each
$(\text{subject},\text{class})$ to one mean-pooled point leaves a handful of high-variance samples on
which a linear probe's post-erasure accuracy is unstable and upward-biased.

Two practical consequences follow. First, single-trial erasure---$n\gg p$, grouped by subject---is the
test that should be reported when claiming an FM is identity-trapped. Second, identity dominance and
per-trial task skill are dissociable: \reve{} is far more subject-dominated than the control yet not
more trap-prone, and a plain SVD beats it per trial on several clean ERP contrasts while the FM wins on
weaker datasets. The honest framing is therefore not ``FMs are trapped'' but ``the pooled erasure
metric manufactures a trap signal that the single-trial metric removes.''

\paragraph{Limitations.} The audit is linear (LEACE + linear probe), as in the original; nonlinear
identity leakage is out of scope. Nine MOABB datasets failed for external reasons (dead hosts, corrupt
archives, or empty class sets after the binary filter). \reve{} subject fraction (a variance share) and
the \tcm{} subject-probe balanced accuracy are related but not identical identity metrics, so we compare
their directions rather than their magnitudes.

\section{Conclusion}
Across the full MOABB benchmark---three paradigms, a foundation model and a classical control, pooled
and per-trial erasure---the identity trap behaves as a pooling artifact: pervasive under the published
pooled metric (including a faithful reproduction of the paper's own MI trap), near-absent
(\fmPtTrapTotal{}/\fmPtNTotal{}) under the honest per-trial metric, and present in a near-identity-free
control. We recommend per-trial, subject-grouped erasure as the standard for identity-trap claims.

\begin{thebibliography}{9}
\bibitem[Belrose et al.(2023)]{belrose2023leace} N.~Belrose et al. LEACE: Perfect linear concept
erasure in closed form. \emph{NeurIPS}, 2023.
\bibitem[\fmscope{}]{fmscope} The Identity Trap in EEG Foundation Models. Preprint, 2026.
\bibitem[Hurlbert(1984)]{hurlbert1984} S.~H.~Hurlbert. Pseudoreplication and the design of ecological
field experiments. \emph{Ecological Monographs}, 54(2):187--211, 1984.
\bibitem[Jayaram \& Barachant(2018)]{jayaram2018moabb} V.~Jayaram and A.~Barachant. MOABB: trustworthy
algorithm benchmarking for BCIs. \emph{J. Neural Eng.}, 2018.
\bibitem[Ravfogel et al.(2020)]{ravfogel2020inlp} S.~Ravfogel, Y.~Elazar, H.~Gonen, M.~Twiton, and
Y.~Goldberg. Null it out: Guarding protected attributes by iterative nullspace projection.
\emph{ACL}, 2020.
\bibitem[Ravfogel et al.(2022)]{ravfogel2022rlace} S.~Ravfogel, M.~Twiton, Y.~Goldberg, and
R.~Cotterell. Linear adversarial concept erasure. \emph{ICML}, 2022.
\bibitem[Wang et al.(2000)]{wang2000tcm} K.~Wang, H.~Begleiter, and B.~Porjesz. Trilinear modeling of
ERP topography (Topographic Component Model). 2000.
\end{thebibliography}

\end{document}
